% !TeX root = ../papers/nc-p3.tex
\section{固定曲線上の長さ5の挿入行列}
\label{sec:fixed-quintic-locus}
\begin{proposition}\label{prop:fixed-quintic-slice}
倍率 $(3,2)$ において両固定曲線上の許容窓の長さ4の次元が全始点で35ならば、
長さ5の次元は全始点で56である。従ってこの倍率で長さ5に初めて現れる障害はない。
\end{proposition}
\begin{proof}
始点 $i$ からの長さ5の道について、外側の二対で関係を割った後の空間は
$Q_i\otimes V_{i+2}\otimes Q_{i+3}$ である。左と右に残る二次関係の挿入は
\[
 (V_i\otimes R_{i+1}\otimes Q_{i+3})\oplus
 (Q_i\otimes R_{i+2}\otimes V_{i+4})
 \longrightarrow Q_i\otimes V_{i+2}\otimes Q_{i+3}.
\]
この余核は長さ5の道空間そのものである。文字ごとの余域次元は100、定義域次元は120となる。
実際に外側二対の文字を選ぶと中央の一文字が決まり、余域の次元は $10^2$。
片方の挿入では関係と対の文字を選べば残る一文字は一意で、次元は $6\cdot10$。
正則一般点では同時文字交換が奇数長の下位ビットを交換する。
さらにKoszul列の群環表示を長さ5まで展開すると、長さ5の文字は四元群の正則表現の14倍である。
実際、奇数次数の各項は正則表現 $v$ を因子に持ち、総次元は56である。
よって目標階数は全て86で、各文字の商次元14は特殊化での下界になる。

\path{computations/verify_fixed_quintic_locus.py} は正負と始点0,1と各文字次数について、
上の行列の二つの $86\times86$ 小行列式を有理数体上で計算する。
始点2,3は二位置移動で戻せる。$b=1$ とし、前節の $(e,f)=(Nu,Du)$ を使う。
各小行列式の既約因子を $G$ との終結式で第一座標に写し、両小行列式に
共通する第一座標の既約因子ごとに曲線上の同時零点を計算する。
$\det A=0$ と既知の全根が不適格な第一座標を除き、
正の側の反対行列式と 可逆表示の非許容条件は曲線上で飽和する。
$u$ の無限遠や $ND=0$ および $a=0$ の上の許容点は、
前の端点計算の長さ5と二位置移動で検査済みである。

負の側の残余は空である。正の側では各始点の文字0,1の検算において、
飽和後に第一座標の既約八次式の根が二組残る。いずれも消去イデアルは
$(P(a),u-U(a))$ の形で、$P$ は既約、$U$ は有理係数である。
各数体について素数1009で $P$ の素体根を一つ取り、$U$ の分母が単元で
実際の挿入行列の階数が86であることを確認する。
従って小行列式の一つはその数体で非零であり、全ての共役点で目標階数に達する。
この還元は既約因子ごとの判定で、素体の点の列挙だけから他の幾何点を除外していない。
\end{proof}
